% ============================================================
% PREÂMBULO DO MODELO DE TESE UFPI
% Objetivo: manter o arquivo organizado, comentado e fácil de editar.
% Observação: este modelo foi pensado para textos matemáticos em português.
% ============================================================

% ------------------------------------------------------------
% IDIOMA E CODIFICAÇÃO
% ------------------------------------------------------------
% babel: traduz automaticamente nomes como "Chapter", "Contents", etc.
\usepackage[brazil]{babel}

% inputenc: permite digitar diretamente com acentos no arquivo-fonte.
\usepackage[utf8]{inputenc}

% fontenc: melhora a saída PDF com palavras acentuadas e hifenização.
\usepackage[T1]{fontenc}

% ------------------------------------------------------------
% MATEMÁTICA
% ------------------------------------------------------------
% amsmath: recursos centrais para fórmulas e ambientes matemáticos.
\usepackage{amsmath}

% amssymb: símbolos matemáticos adicionais.
\usepackage{amssymb}

% amsthm: ambientes de teorema, lema, proposição, definição etc.
\usepackage{amsthm}

% amsfonts: fontes matemáticas adicionais.
\usepackage{amsfonts}

% mathrsfs: caligrafia matemática alternativa, útil para alguns símbolos.
\usepackage{mathrsfs}

% dsfont: permite usar, por exemplo, \mathds{R}, \mathds{N}, etc.
\usepackage{dsfont}

% eulervm: altera o visual da matemática para a família Euler.
\usepackage{eulervm}

% ------------------------------------------------------------
% ENUMERAÇÃO E FORMATAÇÃO GERAL
% ------------------------------------------------------------
% enumerate: controle simples de listas enumeradas.
\usepackage{enumerate}

% geometry: ajusta margens e espaçamentos de página.
\usepackage{geometry}

% fancyhdr: personaliza cabeçalhos e numeração de páginas.
\usepackage{fancyhdr}

% graphicx: insere figuras e logotipos.
\usepackage{graphicx}

%referencias clicáveis.
\usepackage{hyperref}

%notas (opcional) - use \todo{revisar esta definição}  ou \todo[inline]{melhorar esta explicação}.
\setlength{\marginparwidth}{2cm}
\usepackage{todonotes}

%mostra todos os labels exceto referencias e nem citações.
\usepackage[notcite]{showkeys}

% ------------------------------------------------------------
% GEOMETRIA DA PÁGINA
% ------------------------------------------------------------
\geometry{
  a4paper,
  tmargin=2.5cm,
  bmargin=2.5cm,
  lmargin=3cm,
  rmargin=2cm,
  headsep=5mm,
  footskip=1cm
}

% Espaçamento entre linhas.
\renewcommand{\baselinestretch}{1.5}

% ------------------------------------------------------------
% CABEÇALHO E RODAPÉ
% ------------------------------------------------------------
\pagestyle{fancy}
\fancyhead[L]{{\bfseries \nouppercase \leftmark}}
\fancyhead[R]{\thepage}
\fancyfoot{}
\setlength{\headheight}{15pt}

% ------------------------------------------------------------
% AMBIENTES DE TEOREMA
% Padronização recomendada: usar nomes completos no código.
% Mantivemos alguns apelidos antigos por compatibilidade com arquivos já escritos.
% ------------------------------------------------------------

% Estilo em itálico para resultados.
\theoremstyle{plain}
\newtheorem{teorema}{Teorema}[chapter]
\newtheorem{lema}[teorema]{Lema}
\newtheorem{proposicao}[teorema]{Proposição}
\newtheorem{corolario}[teorema]{Corolário}

% Estilo normal para definições e exemplos.
\theoremstyle{definition}
\newtheorem{definicao}[teorema]{Definição}
\newtheorem{exemplo}[teorema]{Exemplo}

% Estilo normal para observações.
\theoremstyle{remark}
\newtheorem{observacao}[teorema]{Observação}

% ------------------------------------------------------------
% COMANDOS MATEMÁTICOS PADRONIZADOS
% ------------------------------------------------------------

% Conjuntos numéricos.
\newcommand{\R}{\mathds{R}}
\newcommand{\N}{\mathds{N}}
\newcommand{\Z}{\mathds{Z}}
\newcommand{\Q}{\mathds{Q}}
\newcommand{\C}{\mathds{C}}

% Letras caligráficas recorrentes.
% Você pode adicionar sua Letra aqui!
\newcommand{\Ac}{\mathcal{A}}
\newcommand{\Cc}{\mathcal{C}}
\newcommand{\Ec}{\mathcal{E}}
\newcommand{\Lc}{\mathcal{L}}
\newcommand{\Nc}{\mathcal{N}}
\newcommand{\Sc}{\mathcal{S}}

% Notações recorrentes já usadas no projeto.
\newcommand{\sig}{\Sigma^n}
\newcommand{\sigg}{\Sigma}
\newcommand{\Ite}[1]{\noindent (#1)}
\newcommand{\denote}{{\noindent\bfseries Notação.\,\,}}
\newcommand{\remarks}{{\noindent\bfseries Observações.\,\,}}
\newcommand{\claim}{{\noindent\bfseries Afirmação.\,\,}}

% ------------------------------------------------------------
% OPERADORES MATEMÁTICOS
% Recomenda-se \DeclareMathOperator para operadores, pois ajusta o espaçamento.
% ------------------------------------------------------------
\DeclareMathOperator{\Int}{int}
\DeclareMathOperator{\supp}{supp}
\DeclareMathOperator{\dist}{dist}
\DeclareMathOperator{\diam}{diam}
\DeclareMathOperator{\osc}{osc}
\DeclareMathOperator{\vol}{vol}
\DeclareMathOperator{\Area}{Area}
\DeclareMathOperator{\grad}{grad}
\DeclareMathOperator{\diver}{div}
\DeclareMathOperator{\Hess}{Hess}
\DeclareMathOperator{\Ric}{Ric}
\DeclareMathOperator{\Rm}{Rm}
\DeclareMathOperator{\Scal}{Scal}

% ------------------------------------------------------------
% AJUSTES ÚTEIS PARA O MODELO
% ------------------------------------------------------------
% Número da equação acompanhando o capítulo: (1.1), (1.2), ...
\numberwithin{equation}{chapter}

% Nomes automáticos em português para referências textuais, se desejado no futuro.
% Exemplo de uso possível com hyperref/cleveref, caso venha a ser incorporado.

